% reformulation.tex --- section fragment.
% Intended to replace the current Section 2 of main.tex via % reformulation.tex --- section fragment.
% Intended to replace the current Section 2 of main.tex via % reformulation.tex --- section fragment.
% Intended to replace the current Section 2 of main.tex via % reformulation.tex --- section fragment.
% Intended to replace the current Section 2 of main.tex via \input{reformulation}.
% Inherits all macros from main.tex (\layer, \Ncut, \vol, \cut, \sign, \argmin,
% \argmax, definition / lemma / proposition theorem envs). Flow operators from
% measure.tex are written as \mathrm{...} to avoid coupling to whether the
% \DeclareMathOperator versions are present in the preamble.

\section{Problem Formulation}\label{sec:setup}

A summary of an attribution graph is determined by a pruned subgraph, a
partition of the surviving nodes into supernodes, and the supernode-level
adjacency that the pair induces. We state the problem in three parts.
Section~\ref{sec:input} fixes the input, Section~\ref{sec:output} fixes the
output and the induced summarization graph, and
Section~\ref{sec:objective} states the objective as the sum of three
losses---coherence, conservation, and complexity---together with
structural well-formedness constraints. Section~\ref{sec:attr-instance}
then commits to a concrete closed form for each loss term. We keep the
problem statement and its instantiation separable from any
\emph{solver}: $K$-means on stacked edge profiles, spectral clustering,
greedy modularity, and the block-coordinate strategy of
Section~\ref{sec:solution} may minimize different surrogates internally,
but every solver is scored on the same $L$ defined here.

\subsection{Input}\label{sec:input}

A directed acyclic attribution graph $G = (V, E, W)$ has vertex partition
$V = V_{\text{emb}} \cup V_{\text{mid}} \cup V_{\text{logit}}$ into
embedding, mid-graph, and logit nodes; edge set $E$; and signed weight
matrix $W \in \mathbb{R}^{N \times N}$ with $N = |V|$. We use the
convention that $W_{uv}$ is the weight on the directed edge $u \to v$,
so $W$ is sparse and strictly upper-triangular in any topological
ordering of $V$. Each node carries a layer index $\layer(\cdot)$;
embeddings sit at layer $0$ and logits at the final layer. Two side data
accompany $G$: a token-weight vector $\tau \in \mathbb{R}^{|V_{\text{emb}}|}$
on the source set (normalized SHAP values, $\sum_e \tau_e = 1$) and a
logit-weight vector $\ell \in \mathbb{R}^{|V_{\text{logit}}|}$ on the
sink set, identifying which inputs and outputs the summary should
explain.

\paragraph{Mid-graph nodes.}
In a cross-layer transcoder attribution graph, $V_{\text{mid}}$
originally contains both \emph{feature} nodes (interpretable
transcoder features) and \emph{reconstruction-error} nodes (the
residual that the transcoder fails to capture at each layer/position).
Error nodes are not interpretable as concept-bearing units, so we
restrict $V_{\text{mid}}$ to feature nodes throughout: all subsequent
references to ``members'' of a supernode, to ``role'' vectors, and to
the partition $\varphi$ apply only to feature nodes.

\subsection{Output}\label{sec:output}

A \emph{summary} is a triple $(V', E', \varphi)$:
\begin{itemize}[leftmargin=*,topsep=2pt,itemsep=1pt]
\item a pruned subgraph $V' \subseteq V$, $E' \subseteq E \cap (V' \times V')$,
with $V_{\text{emb}} \cup V_{\text{logit}} \subseteq V'$ forced;
\item a partition $\varphi : V' \to \pi(\varphi)$ assigning each surviving
node to exactly one \emph{supernode} $S \in \pi(\varphi)$, with every
$v \in V_{\text{emb}} \cup V_{\text{logit}}$ a singleton supernode.
\end{itemize}
The summary induces the \emph{summarization graph}
$G_{SN} = (\pi(\varphi),\, E^{SN},\, W^{SN})$ whose nodes are the
supernodes and whose adjacency aggregates the surviving weights:
\begin{equation}
W^{SN}_{ST} \;=\; \sum_{(u, v) \in E'} W_{uv}\,
\mathbf{1}[u \in S]\,\mathbf{1}[v \in T],
\label{eq:wsn}
\end{equation}
with $E^{SN} = \{(S, T) : W^{SN}_{ST} \ne 0,\; S \ne T\}$. The
summarization graph is the artifact the user reads; $(V', E', \varphi)$
are the decision variables that produce it.

\subsection{Objective}\label{sec:objective}

\paragraph{Three losses.}
A summary is good to the extent that (i) the members of each supernode
share a mechanistic role, (ii) $G_{SN}$ preserves the causal mass that
flows through $G$ from $V_{\text{emb}}$ to $V_{\text{logit}}$, and (iii)
the summary is small enough to read. We assemble these into a single
scalar objective:
\begin{equation}
\min_{(V', E', \varphi)}\;
L(V', E', \varphi)
\;=\;
L_{\text{coh}} \;+\; L_{\text{cons}} \;+\; L_{\text{cplx}}.
\label{eq:Lsum}
\end{equation}
Each term is named by what it measures, not by how it is measured:
\begin{description}[leftmargin=2.0em,style=nextline,topsep=2pt,itemsep=2pt]
\item[$L_{\text{coh}}(\varphi;\, V', E')$.] Within-supernode disagreement
on mechanistic role.
\item[$L_{\text{cons}}(V', E', \varphi;\, G)$.] Causal mass lost when
$G$ is replaced by $G_{SN}$.
\item[$L_{\text{cplx}}(V', E', \varphi)$.] Description length of the
summary.
\end{description}
The closed form of each term---and the multipliers that balance them---is
a property of a \emph{solver}, not of the problem.
Section~\ref{sec:attr-instance} fixes a concrete instantiation for
transcoder attribution graphs; Section~\ref{sec:eval} explains how
candidate solutions to~\eqref{eq:Lsum} are scored.

\paragraph{Structural constraints.}
Independent of the loss, a summary must be structurally well-formed:
\begin{description}[leftmargin=2.0em,style=nextline,topsep=2pt,itemsep=2pt]
\item[(F1) DAG integrity.] $G_{SN}$ is acyclic.
\item[(F2) Partition well-formedness.] $\varphi$ is a function (each
$v \in V'$ in exactly one supernode); embedding and logit nodes are
singletons.
\end{description}
(F1) and (F2) are properties any valid summary must satisfy, regardless of
which loss instantiation or solver is in play. Soft surrogates
(Lagrangian relaxations, penalty terms, dual variables) are introduced
only by specific solvers; the problem statement treats them as hard.

\section{Instantiation for the Transcoder Attribution Graph}\label{sec:attr-instance}

The losses of Section~\ref{sec:objective} are abstract. We now define each
one concretely for the transcoder attribution graph, where feature nodes
correspond to transcoder features and edge weights $W_{uv}$ encode the
linear causal contribution of $u$ to $v$ under the cross-layer
transcoder linearization. The point of this section is to pin down
\emph{what} is being measured; the survey of solvers and their surrogate
forms is the subject of Section~\ref{sec:solution}.

\subsection{Coherence}\label{sec:attr-coh}

The \emph{role} of a feature $u \in V'_{\text{mid}}$ is its pair of
weighted outgoing and incoming edge profiles, with each coordinate
scaled by the causal importance of its endpoint:
\begin{align}
(v_u^{\text{out}})_k &= W_{uk}\,\sqrt{\mathrm{Inf}_k} \nonumber \\
(v_u^{\text{in}})_k  &= W_{ku}\,\sqrt{\mathrm{Rel}_k} \nonumber \\
r(u) &= \bigl[\,v_u^{\text{out}}\,;\, v_u^{\text{in}}\,\bigr] \in \mathbb{R}^{2N}.
\label{eq:role}
\end{align}
Here $\mathrm{Inf}_k$ and $\mathrm{Rel}_k$ are the per-node influence
and relevance scores defined in Appendix~\ref{app:flow} via the Neumann
series on $G$. Outgoing coordinates are weighted by the
\emph{downstream} importance of their target ($\sqrt{\mathrm{Inf}_k}$),
incoming coordinates by the \emph{upstream} importance of their source
($\sqrt{\mathrm{Rel}_k}$), so the role vector concentrates on the edge
dimensions that actually carry causal mass. This is exactly the role
vector computed by every clusterer in our evaluation (the
\texttt{compute\_phi\_vectors} construction).

We commit to a single closed form for the coherence loss: the
within-supernode sum of squared Euclidean distances of each member's
role to the supernode centroid,
\begin{equation}
L_{\text{coh}}(\varphi) \;=\;
\sum_{S \in \pi(\varphi)}\; \sum_{u \in S}
\;\bigl\lVert r(u) - \bar r(S) \bigr\rVert_2^2,
\label{eq:Lcoh}
\end{equation}
where $\bar r(S) = \frac{1}{|S|} \sum_{u \in S} r(u)$ is the supernode
centroid in role space.
Eq.~\eqref{eq:Lcoh} is the $K$-means within-cluster sum of squares on
the Inf/Rel-weighted edge profiles. It is fixed once and reported
uniformly across all solvers: a spectral clusterer that internally
minimizes a normalized cut on a reweighted affinity, a modularity
clusterer that maximizes a degree-corrected score, and a direct
$K$-means clusterer all yield a partition $\varphi$, and each is
scored on the same $L_{\text{coh}}(\varphi)$ above. The internal
surrogate may differ; the reported number does not.

\paragraph{Why this choice.}
The Inf/Rel reweighting matches the role vector actually computed by
the methods we evaluate, and the squared-Euclidean form makes
$L_{\text{coh}}$ a well-defined function of the partition regardless of
which surrogate any solver used internally. The price is that
$\mathrm{Rel}$ inherits its dependence on the token-weight vector
$\tau$, currently set to SHAP values---the same fairness limitation we
flag for $L_{\text{cons}}$ (Section~\ref{sec:attr-cons}). A
steering-based replacement---defining $r(u)$ as the per-feature logit
perturbation $\Delta_u$ under intervention rather than the structural
edge profile, and reusing Eq.~\eqref{eq:Lcoh}---would directly measure
whether grouped features behave alike causally and is the natural
future iteration.

\subsection{Conservation}\label{sec:attr-cons}

Edges in an attribution graph carry signed linear contributions;
composing them along directed paths gives the causal mass that reaches
$V_{\text{logit}}$ from $V_{\text{emb}}$. Conservation is measured in
the same units. Working with magnitudes $|W|$, define the outgoing
transition matrix
\begin{equation}
P^{\text{out}}_{uv}
\;=\;
\frac{|W_{uv}|}{\sum_{w : (u, w) \in E'} |W_{uw}|}
\label{eq:Pout-reform}
\end{equation}
for each non-sink $u \in V'$, and the source vector $F^{\text{src}}$
equal to $\tau$ on $V_{\text{emb}}$ and zero elsewhere. The node-level
flow $\mathrm{Flow}(u) = [F^{\text{src}\top}(I - P^{\text{out}})^{-1}]_u$
is well-defined on a DAG and satisfies a per-node conservation law plus
the end-to-end identity $\sum_{\ell \in V_{\text{logit}}} \mathrm{Flow}(\ell)
= 1$ on the original (unpruned) graph. Aggregating to supernodes,
$\mathrm{Flow}^{SN}(S) = \sum_{u \in S} \mathrm{Flow}(u)$, similarly
balances at every interior supernode of $G_{SN}$. Full statements,
proofs, and the supernode lift are deferred to the flow note that
accompanies this paper.

A summary loses causal mass in two distinct ways.

\paragraph{Pruning loss.}
Edges and nodes removed in the cut to $(V', E')$ carry flow that no
longer reaches a logit. The fraction of source mass that survives
pruning is the \emph{flow completeness}
\begin{equation}
\mathcal{C}(V', E')
\;=\;
\sum_{\ell \in V_{\text{logit}}} \mathrm{Flow}_{V'}(\ell)
\;\in\; [0, 1],
\label{eq:Ccal}
\end{equation}
where $\mathrm{Flow}_{V'}$ denotes flow computed on the pruned graph
with $\tau$ restricted and renormalized. $\mathcal{C}(V', E')$ depends
only on $(V', E')$ and is partition-invariant.

\paragraph{Aggregation loss.}
Contracting features into supernodes destroys edge weight in two ways.
Edges between members of the same supernode become self-loops and are
dropped from $G_{SN}$. Edges between members of different supernodes
sum with sign in $W^{SN}_{ST} = \sum_{u \in S, v \in T} W_{uv}$, so
opposing-sign edges from $S$ to $T$ can partially or fully cancel.
Both forms of loss are visible directly on the edge magnitudes:
\begin{equation}
D_{\text{agg}}(\varphi)
\;=\;
1 \;-\;
\frac{\sum_{(S, T) \in E^{SN}} \lvert W^{SN}_{ST} \rvert}
     {\sum_{(u, v) \in E'} \lvert W_{uv} \rvert}.
\label{eq:Dagg}
\end{equation}
The numerator is the total magnitude of inter-supernode weight retained
in $G_{SN}$; the denominator is the total edge magnitude in the pruned
graph. $D_{\text{agg}}(\varphi) \in [0, 1]$, with $0$ iff every pruned
edge is preserved as a same-sign contribution to a distinct
inter-supernode block, and $1$ iff every pruned edge is absorbed (for
example $|\pi(\varphi)| = 1$, all mid-graph features in one supernode).

\paragraph{Conservation loss.}
Combining the two,
\begin{equation}
L_{\text{cons}}(V', E', \varphi)
\;=\;
\underbrace{\bigl(1 - \mathcal{C}(V', E')\bigr)}_{\text{prune loss}}
\;+\;
\underbrace{D_{\text{agg}}(\varphi)}_{\text{aggregation loss}}.
\label{eq:Lcons}
\end{equation}
The first term is the price of pruning and depends on $(V', E')$ and on
the source vector $\tau$; the second is the price of grouping and
depends only on $\varphi$ given $(V', E')$. No similarity matrix,
signature reweighting, or spectral relaxation appears in either
definition.

\paragraph{Fairness caveat on $\tau$.}
$D_{\text{agg}}$ is structural: it depends only on the signed edge
weights $W$ and the partition $\varphi$, with no role for $\tau$ or
$\ell$, so the aggregation half of $L_{\text{cons}}$ is fairness-clean. The
prune-loss half $1 - \mathcal{C}$ is not: it is precisely the
fraction of $\tau$-injected mass that fails to reach a logit, which
makes $\mathcal{C}(V', E')$ equal to the relevance-reachability of
$V_{\text{logit}}$. Our current pipeline sets $\tau$ to SHAP values,
so the prune loss inherits the bias of an external attribution
method. We treat $\tau$ as a free parameter of the problem
statement---any nonnegative distribution on $V_{\text{emb}}$ with
$\sum_e \tau_e = 1$ is admissible---and we flag a principled fairness
fix as open work: candidate alternatives include uniform $\tau =
1/|V_{\text{emb}}|$, a backward influence-side flow rooted at
$V_{\text{logit}}$ with logit weights $\ell$, and a two-sided
symmetrization combining forward and backward completeness.

\subsection{Complexity}\label{sec:attr-cplx}

A reader follows the summary by tracing causal paths from input tokens
to output logits. The relevant complexity is therefore the typical
distance from $V_{\text{emb}}$ to $V_{\text{logit}}$ in the supernode
graph, measured in supernode hops:
\begin{equation}
L_{\text{cplx}}(V', E', \varphi)
\;=\;
\frac{1}{|\mathcal{P}|}
\sum_{(e, \ell) \in \mathcal{P}}
d_{G_{SN}}(e, \ell),
\label{eq:Lcplx}
\end{equation}
where $\mathcal{P} = \{(e, \ell) \in V_{\text{emb}} \times V_{\text{logit}} : \ell \text{ reachable from } e \text{ in } G_{SN}\}$
and $d_{G_{SN}}(e, \ell)$ is the shortest-path length between $e$ and
$\ell$ in $G_{SN}$, counted in supernode edges. We set $L_{\text{cplx}}
= 0$ by convention when $\mathcal{P} = \emptyset$ (a degenerate
disconnected summary).

Eq.~\eqref{eq:Lcplx} is the plain average shortest-path length over
embedding-logit pairs; no flow weighting and no $\tau$ enters its
definition, so $L_{\text{cplx}}$ is fairness-clean in the sense of the
caveat in Section~\ref{sec:attr-cons}. Collapsing all mid-graph
features into a single supernode trivially hits the floor
$L_{\text{cplx}} = 2$ (one hop into the collapsed supernode, one hop
out), so unbounded grouping is checked by $L_{\text{cons}}$ through
$D_{\text{agg}}$ rather than rewarded here.

\section{Evaluation Methodology}\label{sec:eval}

A candidate solution to~\eqref{eq:Lsum} is a triple $(V', E', \varphi)$.
We score it along three axes: (i) the objective itself, decomposed by
term; (ii) the causal behaviour of the resulting summarization graph
under intervention; (iii) external utility for the human reader.
Decomposing in this way lets us compare solvers that minimize different
surrogates of the same problem on a common scale.

\subsection{Objective-aligned metrics}\label{sec:eval-obj}

Because $L$ is concrete (Section~\ref{sec:attr-instance}), every solver
is scored on exactly the same three numbers regardless of which
surrogate it minimized internally.
\begin{itemize}[leftmargin=*,topsep=2pt,itemsep=1pt]
\item \textbf{Coherence.} The realized $L_{\text{coh}}(\varphi)$ from
Eq.~\eqref{eq:Lcoh}: within-supernode sum of squared Euclidean
distances of the Inf/Rel-weighted role vectors to the supernode
centroid, identical to running \texttt{compute\_phi\_vectors} followed
by within-cluster inertia.
\item \textbf{Conservation.} $\mathcal{C}(V', E')$ and
$D_{\text{agg}}(\varphi)$ from Eqs.~\eqref{eq:Ccal}--\eqref{eq:Dagg},
reported separately. The first is partition-invariant given a fixed
pruner and acts as a sanity check on the pruning step; the second is
the per-partition quantity solvers compete on. $\mathcal{C}$ depends on
the choice of $\tau$, which defaults to SHAP values
(Section~\ref{sec:attr-cons}).
\item \textbf{Complexity.} $L_{\text{cplx}}(\varphi)$ from
Eq.~\eqref{eq:Lcplx}: average shortest-path length from
$V_{\text{emb}}$ to $V_{\text{logit}}$ in $G_{SN}$. We also report
$|\pi(\varphi)|$ as a descriptive statistic, though it is not part of
the objective.
\end{itemize}

\subsection{Causal-behaviour metrics}\label{sec:eval-causal}

$L_{\text{cons}}$ as defined in~\eqref{eq:Lcons} is intrinsic to the
graph. We additionally verify that the resulting $G_{SN}$ behaves
causally as the partition claims. The flow-based measures of the
companion flow note apply directly and are graph-intrinsic (no forward
passes through the model are required):
\begin{itemize}[leftmargin=*,topsep=2pt,itemsep=1pt]
\item \textbf{Flow faithfulness correlation}
$\mathrm{FaithCorr}_{\mathrm{Flow}}(\varphi)$: Spearman correlation
between supernode flow $\mathrm{Flow}^{SN}(S)$ and deletion-induced
flow drop $\Delta \mathcal{C}(S)$. A faithful partition ranks supernodes
by flow consistently with their causal contribution to logit-bound
flow.
\item \textbf{Flow-DAUC} and \textbf{Flow-IAUC}: areas under the
deletion and insertion curves over supernodes ranked by flow.
Lower DAUC and higher IAUC indicate that flow-ranked supernodes carry
the causal mass the rank suggests.
\item \textbf{Monotonicity} $\mathrm{Mono}(\varphi)$: fraction of
consecutive ranked supernodes whose deletion drops are non-increasing.
\item \textbf{Sufficiency} $\mathrm{Suff}_k^{\mathrm{Flow}}$ and
\textbf{comprehensiveness} $\mathrm{Comp}_k^{\mathrm{Flow}}$ at fixed
top-$k$: how much sink flow survives keeping or removing only the
top-$k$ supernodes.
\end{itemize}
Each of these scores a partition $\varphi$ using only $(V', E', \varphi)$
and the source vector $\tau$, with no reference to the solver that
produced $\varphi$.

\subsection{External utility metrics}\label{sec:eval-ext}

Two further metrics live outside the optimization problem but are the
practical reason summarization is useful.

\paragraph{Semantic coherence (LLM judge).}
For each supernode $S$, prompt a judge model with the auto-interpretation
of every member feature and ask for a coherence score in $[0, 1]$;
average over supernodes. This measures whether the partition aligns with
human-readable concepts, which $L_{\text{coh}}$ approximates structurally
but does not enforce.

\paragraph{Steering.}
For each supernode $S$ in a held-out evaluation set, scale the
activations of all member features by a common factor at inference time
and measure the resulting shift in the target logit. A partition
supports steering if intervening at the supernode level produces
predictable and concept-aligned logit shifts. This is the operational
test of whether $G_{SN}$ exposes causally controllable handles.

\subsection{Two-stage evaluation}\label{sec:eval-compare}

The pipeline factors into a pruning stage (producing $(V', E')$) and
an aggregation stage (producing $\varphi$), each with its own solver.
We evaluate the two stages independently. Pruning is scored through
the prune-loss component $1 - \mathcal{C}(V', E')$ alone, on a
candidate $(V', E')$ basis, and is reported separately. The protocol
below is the aggregation comparison, with the pruner held fixed.

\paragraph{Fixed pruner.}
A single pruning algorithm produces the $(V', E')$ that feeds every
aggregation solver in the comparison. $\mathcal{C}(V', E')$ is
therefore identical across rows, the prune-loss contribution to
$L_{\text{cons}}$ is the same constant in every row, and all
differences come from $D_{\text{agg}}$, $L_{\text{coh}}$, and
$L_{\text{cplx}}$.

\paragraph{Per-solver best-configuration reporting.}
Each aggregation solver has its own hyperparameter space; we sweep
exhaustively and report \emph{only the best configuration per solver}
in the comparison table. The hyperparameter spaces are:
\begin{itemize}[leftmargin=*,topsep=2pt,itemsep=1pt]
\item \textbf{ASG-spectral (ours)}: number of clusters $K$, mean
operator $F_{\text{mean}} \in \{\text{arith}, \text{harm}, \text{geo}\}$,
layer-decay rate $\lambda \in [0, 1]$.
\item \textbf{ASG-agglomerative (ours)}: same $K, F_{\text{mean}}, \lambda$,
with agglomerative linkage replacing the spectral relaxation.
\item \textbf{Greedy modularity}: target community count $K$ on the
symmetrized weighted adjacency.
\item \textbf{Spectral-RBF}: $K$, RBF kernel over the Inf/Rel-weighted
role vectors $r(u)$.
\item \textbf{$K$-means}: $K$ on the same role vectors.
\end{itemize}
``Best'' is decided by a single selection criterion fixed in advance
and shared across all solvers (the experiments section gives the exact
form). $K$ ranges over an eigengap-bounded candidate set common to
every solver. Different solvers may select different $K^\star$, which
is permitted: matching $|\pi(\varphi)|$ across solvers would force
sub-optimal configurations on methods whose natural granularity
differs, and the metrics already account for granularity via
$L_{\text{cplx}}$.

\paragraph{Reporting and dominance.}
For each solver's best configuration we report all metric families
from Sections~\ref{sec:eval-obj}--\ref{sec:eval-ext}. A solver
dominates if it is Pareto-better on the objective-aligned metrics
($L_{\text{coh}}, D_{\text{agg}}, L_{\text{cplx}}$) and at least
matches on causal-behaviour and external metrics. Because every
objective-aligned metric is intrinsic to the supernode graph---no
method-private similarity matrix or spectral surrogate enters its
definition---a solver that does well here is doing well at the stated
problem rather than at its own internal proxy.
.
% Inherits all macros from main.tex (\layer, \Ncut, \vol, \cut, \sign, \argmin,
% \argmax, definition / lemma / proposition theorem envs). Flow operators from
% measure.tex are written as \mathrm{...} to avoid coupling to whether the
% \DeclareMathOperator versions are present in the preamble.

\section{Problem Formulation}\label{sec:setup}

A summary of an attribution graph is determined by a pruned subgraph, a
partition of the surviving nodes into supernodes, and the supernode-level
adjacency that the pair induces. We state the problem in three parts.
Section~\ref{sec:input} fixes the input, Section~\ref{sec:output} fixes the
output and the induced summarization graph, and
Section~\ref{sec:objective} states the objective as the sum of three
losses---coherence, conservation, and complexity---together with
structural well-formedness constraints. Section~\ref{sec:attr-instance}
then commits to a concrete closed form for each loss term. We keep the
problem statement and its instantiation separable from any
\emph{solver}: $K$-means on stacked edge profiles, spectral clustering,
greedy modularity, and the block-coordinate strategy of
Section~\ref{sec:solution} may minimize different surrogates internally,
but every solver is scored on the same $L$ defined here.

\subsection{Input}\label{sec:input}

A directed acyclic attribution graph $G = (V, E, W)$ has vertex partition
$V = V_{\text{emb}} \cup V_{\text{mid}} \cup V_{\text{logit}}$ into
embedding, mid-graph, and logit nodes; edge set $E$; and signed weight
matrix $W \in \mathbb{R}^{N \times N}$ with $N = |V|$. We use the
convention that $W_{uv}$ is the weight on the directed edge $u \to v$,
so $W$ is sparse and strictly upper-triangular in any topological
ordering of $V$. Each node carries a layer index $\layer(\cdot)$;
embeddings sit at layer $0$ and logits at the final layer. Two side data
accompany $G$: a token-weight vector $\tau \in \mathbb{R}^{|V_{\text{emb}}|}$
on the source set (normalized SHAP values, $\sum_e \tau_e = 1$) and a
logit-weight vector $\ell \in \mathbb{R}^{|V_{\text{logit}}|}$ on the
sink set, identifying which inputs and outputs the summary should
explain.

\paragraph{Mid-graph nodes.}
In a cross-layer transcoder attribution graph, $V_{\text{mid}}$
originally contains both \emph{feature} nodes (interpretable
transcoder features) and \emph{reconstruction-error} nodes (the
residual that the transcoder fails to capture at each layer/position).
Error nodes are not interpretable as concept-bearing units, so we
restrict $V_{\text{mid}}$ to feature nodes throughout: all subsequent
references to ``members'' of a supernode, to ``role'' vectors, and to
the partition $\varphi$ apply only to feature nodes.

\subsection{Output}\label{sec:output}

A \emph{summary} is a triple $(V', E', \varphi)$:
\begin{itemize}[leftmargin=*,topsep=2pt,itemsep=1pt]
\item a pruned subgraph $V' \subseteq V$, $E' \subseteq E \cap (V' \times V')$,
with $V_{\text{emb}} \cup V_{\text{logit}} \subseteq V'$ forced;
\item a partition $\varphi : V' \to \pi(\varphi)$ assigning each surviving
node to exactly one \emph{supernode} $S \in \pi(\varphi)$, with every
$v \in V_{\text{emb}} \cup V_{\text{logit}}$ a singleton supernode.
\end{itemize}
The summary induces the \emph{summarization graph}
$G_{SN} = (\pi(\varphi),\, E^{SN},\, W^{SN})$ whose nodes are the
supernodes and whose adjacency aggregates the surviving weights:
\begin{equation}
W^{SN}_{ST} \;=\; \sum_{(u, v) \in E'} W_{uv}\,
\mathbf{1}[u \in S]\,\mathbf{1}[v \in T],
\label{eq:wsn}
\end{equation}
with $E^{SN} = \{(S, T) : W^{SN}_{ST} \ne 0,\; S \ne T\}$. The
summarization graph is the artifact the user reads; $(V', E', \varphi)$
are the decision variables that produce it.

\subsection{Objective}\label{sec:objective}

\paragraph{Three losses.}
A summary is good to the extent that (i) the members of each supernode
share a mechanistic role, (ii) $G_{SN}$ preserves the causal mass that
flows through $G$ from $V_{\text{emb}}$ to $V_{\text{logit}}$, and (iii)
the summary is small enough to read. We assemble these into a single
scalar objective:
\begin{equation}
\min_{(V', E', \varphi)}\;
L(V', E', \varphi)
\;=\;
L_{\text{coh}} \;+\; L_{\text{cons}} \;+\; L_{\text{cplx}}.
\label{eq:Lsum}
\end{equation}
Each term is named by what it measures, not by how it is measured:
\begin{description}[leftmargin=2.0em,style=nextline,topsep=2pt,itemsep=2pt]
\item[$L_{\text{coh}}(\varphi;\, V', E')$.] Within-supernode disagreement
on mechanistic role.
\item[$L_{\text{cons}}(V', E', \varphi;\, G)$.] Causal mass lost when
$G$ is replaced by $G_{SN}$.
\item[$L_{\text{cplx}}(V', E', \varphi)$.] Description length of the
summary.
\end{description}
The closed form of each term---and the multipliers that balance them---is
a property of a \emph{solver}, not of the problem.
Section~\ref{sec:attr-instance} fixes a concrete instantiation for
transcoder attribution graphs; Section~\ref{sec:eval} explains how
candidate solutions to~\eqref{eq:Lsum} are scored.

\paragraph{Structural constraints.}
Independent of the loss, a summary must be structurally well-formed:
\begin{description}[leftmargin=2.0em,style=nextline,topsep=2pt,itemsep=2pt]
\item[(F1) DAG integrity.] $G_{SN}$ is acyclic.
\item[(F2) Partition well-formedness.] $\varphi$ is a function (each
$v \in V'$ in exactly one supernode); embedding and logit nodes are
singletons.
\end{description}
(F1) and (F2) are properties any valid summary must satisfy, regardless of
which loss instantiation or solver is in play. Soft surrogates
(Lagrangian relaxations, penalty terms, dual variables) are introduced
only by specific solvers; the problem statement treats them as hard.

\section{Instantiation for the Transcoder Attribution Graph}\label{sec:attr-instance}

The losses of Section~\ref{sec:objective} are abstract. We now define each
one concretely for the transcoder attribution graph, where feature nodes
correspond to transcoder features and edge weights $W_{uv}$ encode the
linear causal contribution of $u$ to $v$ under the cross-layer
transcoder linearization. The point of this section is to pin down
\emph{what} is being measured; the survey of solvers and their surrogate
forms is the subject of Section~\ref{sec:solution}.

\subsection{Coherence}\label{sec:attr-coh}

The \emph{role} of a feature $u \in V'_{\text{mid}}$ is its pair of
weighted outgoing and incoming edge profiles, with each coordinate
scaled by the causal importance of its endpoint:
\begin{align}
(v_u^{\text{out}})_k &= W_{uk}\,\sqrt{\mathrm{Inf}_k} \nonumber \\
(v_u^{\text{in}})_k  &= W_{ku}\,\sqrt{\mathrm{Rel}_k} \nonumber \\
r(u) &= \bigl[\,v_u^{\text{out}}\,;\, v_u^{\text{in}}\,\bigr] \in \mathbb{R}^{2N}.
\label{eq:role}
\end{align}
Here $\mathrm{Inf}_k$ and $\mathrm{Rel}_k$ are the per-node influence
and relevance scores defined in Appendix~\ref{app:flow} via the Neumann
series on $G$. Outgoing coordinates are weighted by the
\emph{downstream} importance of their target ($\sqrt{\mathrm{Inf}_k}$),
incoming coordinates by the \emph{upstream} importance of their source
($\sqrt{\mathrm{Rel}_k}$), so the role vector concentrates on the edge
dimensions that actually carry causal mass. This is exactly the role
vector computed by every clusterer in our evaluation (the
\texttt{compute\_phi\_vectors} construction).

We commit to a single closed form for the coherence loss: the
within-supernode sum of squared Euclidean distances of each member's
role to the supernode centroid,
\begin{equation}
L_{\text{coh}}(\varphi) \;=\;
\sum_{S \in \pi(\varphi)}\; \sum_{u \in S}
\;\bigl\lVert r(u) - \bar r(S) \bigr\rVert_2^2,
\label{eq:Lcoh}
\end{equation}
where $\bar r(S) = \frac{1}{|S|} \sum_{u \in S} r(u)$ is the supernode
centroid in role space.
Eq.~\eqref{eq:Lcoh} is the $K$-means within-cluster sum of squares on
the Inf/Rel-weighted edge profiles. It is fixed once and reported
uniformly across all solvers: a spectral clusterer that internally
minimizes a normalized cut on a reweighted affinity, a modularity
clusterer that maximizes a degree-corrected score, and a direct
$K$-means clusterer all yield a partition $\varphi$, and each is
scored on the same $L_{\text{coh}}(\varphi)$ above. The internal
surrogate may differ; the reported number does not.

\paragraph{Why this choice.}
The Inf/Rel reweighting matches the role vector actually computed by
the methods we evaluate, and the squared-Euclidean form makes
$L_{\text{coh}}$ a well-defined function of the partition regardless of
which surrogate any solver used internally. The price is that
$\mathrm{Rel}$ inherits its dependence on the token-weight vector
$\tau$, currently set to SHAP values---the same fairness limitation we
flag for $L_{\text{cons}}$ (Section~\ref{sec:attr-cons}). A
steering-based replacement---defining $r(u)$ as the per-feature logit
perturbation $\Delta_u$ under intervention rather than the structural
edge profile, and reusing Eq.~\eqref{eq:Lcoh}---would directly measure
whether grouped features behave alike causally and is the natural
future iteration.

\subsection{Conservation}\label{sec:attr-cons}

Edges in an attribution graph carry signed linear contributions;
composing them along directed paths gives the causal mass that reaches
$V_{\text{logit}}$ from $V_{\text{emb}}$. Conservation is measured in
the same units. Working with magnitudes $|W|$, define the outgoing
transition matrix
\begin{equation}
P^{\text{out}}_{uv}
\;=\;
\frac{|W_{uv}|}{\sum_{w : (u, w) \in E'} |W_{uw}|}
\label{eq:Pout-reform}
\end{equation}
for each non-sink $u \in V'$, and the source vector $F^{\text{src}}$
equal to $\tau$ on $V_{\text{emb}}$ and zero elsewhere. The node-level
flow $\mathrm{Flow}(u) = [F^{\text{src}\top}(I - P^{\text{out}})^{-1}]_u$
is well-defined on a DAG and satisfies a per-node conservation law plus
the end-to-end identity $\sum_{\ell \in V_{\text{logit}}} \mathrm{Flow}(\ell)
= 1$ on the original (unpruned) graph. Aggregating to supernodes,
$\mathrm{Flow}^{SN}(S) = \sum_{u \in S} \mathrm{Flow}(u)$, similarly
balances at every interior supernode of $G_{SN}$. Full statements,
proofs, and the supernode lift are deferred to the flow note that
accompanies this paper.

A summary loses causal mass in two distinct ways.

\paragraph{Pruning loss.}
Edges and nodes removed in the cut to $(V', E')$ carry flow that no
longer reaches a logit. The fraction of source mass that survives
pruning is the \emph{flow completeness}
\begin{equation}
\mathcal{C}(V', E')
\;=\;
\sum_{\ell \in V_{\text{logit}}} \mathrm{Flow}_{V'}(\ell)
\;\in\; [0, 1],
\label{eq:Ccal}
\end{equation}
where $\mathrm{Flow}_{V'}$ denotes flow computed on the pruned graph
with $\tau$ restricted and renormalized. $\mathcal{C}(V', E')$ depends
only on $(V', E')$ and is partition-invariant.

\paragraph{Aggregation loss.}
Contracting features into supernodes destroys edge weight in two ways.
Edges between members of the same supernode become self-loops and are
dropped from $G_{SN}$. Edges between members of different supernodes
sum with sign in $W^{SN}_{ST} = \sum_{u \in S, v \in T} W_{uv}$, so
opposing-sign edges from $S$ to $T$ can partially or fully cancel.
Both forms of loss are visible directly on the edge magnitudes:
\begin{equation}
D_{\text{agg}}(\varphi)
\;=\;
1 \;-\;
\frac{\sum_{(S, T) \in E^{SN}} \lvert W^{SN}_{ST} \rvert}
     {\sum_{(u, v) \in E'} \lvert W_{uv} \rvert}.
\label{eq:Dagg}
\end{equation}
The numerator is the total magnitude of inter-supernode weight retained
in $G_{SN}$; the denominator is the total edge magnitude in the pruned
graph. $D_{\text{agg}}(\varphi) \in [0, 1]$, with $0$ iff every pruned
edge is preserved as a same-sign contribution to a distinct
inter-supernode block, and $1$ iff every pruned edge is absorbed (for
example $|\pi(\varphi)| = 1$, all mid-graph features in one supernode).

\paragraph{Conservation loss.}
Combining the two,
\begin{equation}
L_{\text{cons}}(V', E', \varphi)
\;=\;
\underbrace{\bigl(1 - \mathcal{C}(V', E')\bigr)}_{\text{prune loss}}
\;+\;
\underbrace{D_{\text{agg}}(\varphi)}_{\text{aggregation loss}}.
\label{eq:Lcons}
\end{equation}
The first term is the price of pruning and depends on $(V', E')$ and on
the source vector $\tau$; the second is the price of grouping and
depends only on $\varphi$ given $(V', E')$. No similarity matrix,
signature reweighting, or spectral relaxation appears in either
definition.

\paragraph{Fairness caveat on $\tau$.}
$D_{\text{agg}}$ is structural: it depends only on the signed edge
weights $W$ and the partition $\varphi$, with no role for $\tau$ or
$\ell$, so the aggregation half of $L_{\text{cons}}$ is fairness-clean. The
prune-loss half $1 - \mathcal{C}$ is not: it is precisely the
fraction of $\tau$-injected mass that fails to reach a logit, which
makes $\mathcal{C}(V', E')$ equal to the relevance-reachability of
$V_{\text{logit}}$. Our current pipeline sets $\tau$ to SHAP values,
so the prune loss inherits the bias of an external attribution
method. We treat $\tau$ as a free parameter of the problem
statement---any nonnegative distribution on $V_{\text{emb}}$ with
$\sum_e \tau_e = 1$ is admissible---and we flag a principled fairness
fix as open work: candidate alternatives include uniform $\tau =
1/|V_{\text{emb}}|$, a backward influence-side flow rooted at
$V_{\text{logit}}$ with logit weights $\ell$, and a two-sided
symmetrization combining forward and backward completeness.

\subsection{Complexity}\label{sec:attr-cplx}

A reader follows the summary by tracing causal paths from input tokens
to output logits. The relevant complexity is therefore the typical
distance from $V_{\text{emb}}$ to $V_{\text{logit}}$ in the supernode
graph, measured in supernode hops:
\begin{equation}
L_{\text{cplx}}(V', E', \varphi)
\;=\;
\frac{1}{|\mathcal{P}|}
\sum_{(e, \ell) \in \mathcal{P}}
d_{G_{SN}}(e, \ell),
\label{eq:Lcplx}
\end{equation}
where $\mathcal{P} = \{(e, \ell) \in V_{\text{emb}} \times V_{\text{logit}} : \ell \text{ reachable from } e \text{ in } G_{SN}\}$
and $d_{G_{SN}}(e, \ell)$ is the shortest-path length between $e$ and
$\ell$ in $G_{SN}$, counted in supernode edges. We set $L_{\text{cplx}}
= 0$ by convention when $\mathcal{P} = \emptyset$ (a degenerate
disconnected summary).

Eq.~\eqref{eq:Lcplx} is the plain average shortest-path length over
embedding-logit pairs; no flow weighting and no $\tau$ enters its
definition, so $L_{\text{cplx}}$ is fairness-clean in the sense of the
caveat in Section~\ref{sec:attr-cons}. Collapsing all mid-graph
features into a single supernode trivially hits the floor
$L_{\text{cplx}} = 2$ (one hop into the collapsed supernode, one hop
out), so unbounded grouping is checked by $L_{\text{cons}}$ through
$D_{\text{agg}}$ rather than rewarded here.

\section{Evaluation Methodology}\label{sec:eval}

A candidate solution to~\eqref{eq:Lsum} is a triple $(V', E', \varphi)$.
We score it along three axes: (i) the objective itself, decomposed by
term; (ii) the causal behaviour of the resulting summarization graph
under intervention; (iii) external utility for the human reader.
Decomposing in this way lets us compare solvers that minimize different
surrogates of the same problem on a common scale.

\subsection{Objective-aligned metrics}\label{sec:eval-obj}

Because $L$ is concrete (Section~\ref{sec:attr-instance}), every solver
is scored on exactly the same three numbers regardless of which
surrogate it minimized internally.
\begin{itemize}[leftmargin=*,topsep=2pt,itemsep=1pt]
\item \textbf{Coherence.} The realized $L_{\text{coh}}(\varphi)$ from
Eq.~\eqref{eq:Lcoh}: within-supernode sum of squared Euclidean
distances of the Inf/Rel-weighted role vectors to the supernode
centroid, identical to running \texttt{compute\_phi\_vectors} followed
by within-cluster inertia.
\item \textbf{Conservation.} $\mathcal{C}(V', E')$ and
$D_{\text{agg}}(\varphi)$ from Eqs.~\eqref{eq:Ccal}--\eqref{eq:Dagg},
reported separately. The first is partition-invariant given a fixed
pruner and acts as a sanity check on the pruning step; the second is
the per-partition quantity solvers compete on. $\mathcal{C}$ depends on
the choice of $\tau$, which defaults to SHAP values
(Section~\ref{sec:attr-cons}).
\item \textbf{Complexity.} $L_{\text{cplx}}(\varphi)$ from
Eq.~\eqref{eq:Lcplx}: average shortest-path length from
$V_{\text{emb}}$ to $V_{\text{logit}}$ in $G_{SN}$. We also report
$|\pi(\varphi)|$ as a descriptive statistic, though it is not part of
the objective.
\end{itemize}

\subsection{Causal-behaviour metrics}\label{sec:eval-causal}

$L_{\text{cons}}$ as defined in~\eqref{eq:Lcons} is intrinsic to the
graph. We additionally verify that the resulting $G_{SN}$ behaves
causally as the partition claims. The flow-based measures of the
companion flow note apply directly and are graph-intrinsic (no forward
passes through the model are required):
\begin{itemize}[leftmargin=*,topsep=2pt,itemsep=1pt]
\item \textbf{Flow faithfulness correlation}
$\mathrm{FaithCorr}_{\mathrm{Flow}}(\varphi)$: Spearman correlation
between supernode flow $\mathrm{Flow}^{SN}(S)$ and deletion-induced
flow drop $\Delta \mathcal{C}(S)$. A faithful partition ranks supernodes
by flow consistently with their causal contribution to logit-bound
flow.
\item \textbf{Flow-DAUC} and \textbf{Flow-IAUC}: areas under the
deletion and insertion curves over supernodes ranked by flow.
Lower DAUC and higher IAUC indicate that flow-ranked supernodes carry
the causal mass the rank suggests.
\item \textbf{Monotonicity} $\mathrm{Mono}(\varphi)$: fraction of
consecutive ranked supernodes whose deletion drops are non-increasing.
\item \textbf{Sufficiency} $\mathrm{Suff}_k^{\mathrm{Flow}}$ and
\textbf{comprehensiveness} $\mathrm{Comp}_k^{\mathrm{Flow}}$ at fixed
top-$k$: how much sink flow survives keeping or removing only the
top-$k$ supernodes.
\end{itemize}
Each of these scores a partition $\varphi$ using only $(V', E', \varphi)$
and the source vector $\tau$, with no reference to the solver that
produced $\varphi$.

\subsection{External utility metrics}\label{sec:eval-ext}

Two further metrics live outside the optimization problem but are the
practical reason summarization is useful.

\paragraph{Semantic coherence (LLM judge).}
For each supernode $S$, prompt a judge model with the auto-interpretation
of every member feature and ask for a coherence score in $[0, 1]$;
average over supernodes. This measures whether the partition aligns with
human-readable concepts, which $L_{\text{coh}}$ approximates structurally
but does not enforce.

\paragraph{Steering.}
For each supernode $S$ in a held-out evaluation set, scale the
activations of all member features by a common factor at inference time
and measure the resulting shift in the target logit. A partition
supports steering if intervening at the supernode level produces
predictable and concept-aligned logit shifts. This is the operational
test of whether $G_{SN}$ exposes causally controllable handles.

\subsection{Two-stage evaluation}\label{sec:eval-compare}

The pipeline factors into a pruning stage (producing $(V', E')$) and
an aggregation stage (producing $\varphi$), each with its own solver.
We evaluate the two stages independently. Pruning is scored through
the prune-loss component $1 - \mathcal{C}(V', E')$ alone, on a
candidate $(V', E')$ basis, and is reported separately. The protocol
below is the aggregation comparison, with the pruner held fixed.

\paragraph{Fixed pruner.}
A single pruning algorithm produces the $(V', E')$ that feeds every
aggregation solver in the comparison. $\mathcal{C}(V', E')$ is
therefore identical across rows, the prune-loss contribution to
$L_{\text{cons}}$ is the same constant in every row, and all
differences come from $D_{\text{agg}}$, $L_{\text{coh}}$, and
$L_{\text{cplx}}$.

\paragraph{Per-solver best-configuration reporting.}
Each aggregation solver has its own hyperparameter space; we sweep
exhaustively and report \emph{only the best configuration per solver}
in the comparison table. The hyperparameter spaces are:
\begin{itemize}[leftmargin=*,topsep=2pt,itemsep=1pt]
\item \textbf{ASG-spectral (ours)}: number of clusters $K$, mean
operator $F_{\text{mean}} \in \{\text{arith}, \text{harm}, \text{geo}\}$,
layer-decay rate $\lambda \in [0, 1]$.
\item \textbf{ASG-agglomerative (ours)}: same $K, F_{\text{mean}}, \lambda$,
with agglomerative linkage replacing the spectral relaxation.
\item \textbf{Greedy modularity}: target community count $K$ on the
symmetrized weighted adjacency.
\item \textbf{Spectral-RBF}: $K$, RBF kernel over the Inf/Rel-weighted
role vectors $r(u)$.
\item \textbf{$K$-means}: $K$ on the same role vectors.
\end{itemize}
``Best'' is decided by a single selection criterion fixed in advance
and shared across all solvers (the experiments section gives the exact
form). $K$ ranges over an eigengap-bounded candidate set common to
every solver. Different solvers may select different $K^\star$, which
is permitted: matching $|\pi(\varphi)|$ across solvers would force
sub-optimal configurations on methods whose natural granularity
differs, and the metrics already account for granularity via
$L_{\text{cplx}}$.

\paragraph{Reporting and dominance.}
For each solver's best configuration we report all metric families
from Sections~\ref{sec:eval-obj}--\ref{sec:eval-ext}. A solver
dominates if it is Pareto-better on the objective-aligned metrics
($L_{\text{coh}}, D_{\text{agg}}, L_{\text{cplx}}$) and at least
matches on causal-behaviour and external metrics. Because every
objective-aligned metric is intrinsic to the supernode graph---no
method-private similarity matrix or spectral surrogate enters its
definition---a solver that does well here is doing well at the stated
problem rather than at its own internal proxy.
.
% Inherits all macros from main.tex (\layer, \Ncut, \vol, \cut, \sign, \argmin,
% \argmax, definition / lemma / proposition theorem envs). Flow operators from
% measure.tex are written as \mathrm{...} to avoid coupling to whether the
% \DeclareMathOperator versions are present in the preamble.

\section{Problem Formulation}\label{sec:setup}

A summary of an attribution graph is determined by a pruned subgraph, a
partition of the surviving nodes into supernodes, and the supernode-level
adjacency that the pair induces. We state the problem in three parts.
Section~\ref{sec:input} fixes the input, Section~\ref{sec:output} fixes the
output and the induced summarization graph, and
Section~\ref{sec:objective} states the objective as the sum of three
losses---coherence, conservation, and complexity---together with
structural well-formedness constraints. Section~\ref{sec:attr-instance}
then commits to a concrete closed form for each loss term. We keep the
problem statement and its instantiation separable from any
\emph{solver}: $K$-means on stacked edge profiles, spectral clustering,
greedy modularity, and the block-coordinate strategy of
Section~\ref{sec:solution} may minimize different surrogates internally,
but every solver is scored on the same $L$ defined here.

\subsection{Input}\label{sec:input}

A directed acyclic attribution graph $G = (V, E, W)$ has vertex partition
$V = V_{\text{emb}} \cup V_{\text{mid}} \cup V_{\text{logit}}$ into
embedding, mid-graph, and logit nodes; edge set $E$; and signed weight
matrix $W \in \mathbb{R}^{N \times N}$ with $N = |V|$. We use the
convention that $W_{uv}$ is the weight on the directed edge $u \to v$,
so $W$ is sparse and strictly upper-triangular in any topological
ordering of $V$. Each node carries a layer index $\layer(\cdot)$;
embeddings sit at layer $0$ and logits at the final layer. Two side data
accompany $G$: a token-weight vector $\tau \in \mathbb{R}^{|V_{\text{emb}}|}$
on the source set (normalized SHAP values, $\sum_e \tau_e = 1$) and a
logit-weight vector $\ell \in \mathbb{R}^{|V_{\text{logit}}|}$ on the
sink set, identifying which inputs and outputs the summary should
explain.

\paragraph{Mid-graph nodes.}
In a cross-layer transcoder attribution graph, $V_{\text{mid}}$
originally contains both \emph{feature} nodes (interpretable
transcoder features) and \emph{reconstruction-error} nodes (the
residual that the transcoder fails to capture at each layer/position).
Error nodes are not interpretable as concept-bearing units, so we
restrict $V_{\text{mid}}$ to feature nodes throughout: all subsequent
references to ``members'' of a supernode, to ``role'' vectors, and to
the partition $\varphi$ apply only to feature nodes.

\subsection{Output}\label{sec:output}

A \emph{summary} is a triple $(V', E', \varphi)$:
\begin{itemize}[leftmargin=*,topsep=2pt,itemsep=1pt]
\item a pruned subgraph $V' \subseteq V$, $E' \subseteq E \cap (V' \times V')$,
with $V_{\text{emb}} \cup V_{\text{logit}} \subseteq V'$ forced;
\item a partition $\varphi : V' \to \pi(\varphi)$ assigning each surviving
node to exactly one \emph{supernode} $S \in \pi(\varphi)$, with every
$v \in V_{\text{emb}} \cup V_{\text{logit}}$ a singleton supernode.
\end{itemize}
The summary induces the \emph{summarization graph}
$G_{SN} = (\pi(\varphi),\, E^{SN},\, W^{SN})$ whose nodes are the
supernodes and whose adjacency aggregates the surviving weights:
\begin{equation}
W^{SN}_{ST} \;=\; \sum_{(u, v) \in E'} W_{uv}\,
\mathbf{1}[u \in S]\,\mathbf{1}[v \in T],
\label{eq:wsn}
\end{equation}
with $E^{SN} = \{(S, T) : W^{SN}_{ST} \ne 0,\; S \ne T\}$. The
summarization graph is the artifact the user reads; $(V', E', \varphi)$
are the decision variables that produce it.

\subsection{Objective}\label{sec:objective}

\paragraph{Three losses.}
A summary is good to the extent that (i) the members of each supernode
share a mechanistic role, (ii) $G_{SN}$ preserves the causal mass that
flows through $G$ from $V_{\text{emb}}$ to $V_{\text{logit}}$, and (iii)
the summary is small enough to read. We assemble these into a single
scalar objective:
\begin{equation}
\min_{(V', E', \varphi)}\;
L(V', E', \varphi)
\;=\;
L_{\text{coh}} \;+\; L_{\text{cons}} \;+\; L_{\text{cplx}}.
\label{eq:Lsum}
\end{equation}
Each term is named by what it measures, not by how it is measured:
\begin{description}[leftmargin=2.0em,style=nextline,topsep=2pt,itemsep=2pt]
\item[$L_{\text{coh}}(\varphi;\, V', E')$.] Within-supernode disagreement
on mechanistic role.
\item[$L_{\text{cons}}(V', E', \varphi;\, G)$.] Causal mass lost when
$G$ is replaced by $G_{SN}$.
\item[$L_{\text{cplx}}(V', E', \varphi)$.] Description length of the
summary.
\end{description}
The closed form of each term---and the multipliers that balance them---is
a property of a \emph{solver}, not of the problem.
Section~\ref{sec:attr-instance} fixes a concrete instantiation for
transcoder attribution graphs; Section~\ref{sec:eval} explains how
candidate solutions to~\eqref{eq:Lsum} are scored.

\paragraph{Structural constraints.}
Independent of the loss, a summary must be structurally well-formed:
\begin{description}[leftmargin=2.0em,style=nextline,topsep=2pt,itemsep=2pt]
\item[(F1) DAG integrity.] $G_{SN}$ is acyclic.
\item[(F2) Partition well-formedness.] $\varphi$ is a function (each
$v \in V'$ in exactly one supernode); embedding and logit nodes are
singletons.
\end{description}
(F1) and (F2) are properties any valid summary must satisfy, regardless of
which loss instantiation or solver is in play. Soft surrogates
(Lagrangian relaxations, penalty terms, dual variables) are introduced
only by specific solvers; the problem statement treats them as hard.

\section{Instantiation for the Transcoder Attribution Graph}\label{sec:attr-instance}

The losses of Section~\ref{sec:objective} are abstract. We now define each
one concretely for the transcoder attribution graph, where feature nodes
correspond to transcoder features and edge weights $W_{uv}$ encode the
linear causal contribution of $u$ to $v$ under the cross-layer
transcoder linearization. The point of this section is to pin down
\emph{what} is being measured; the survey of solvers and their surrogate
forms is the subject of Section~\ref{sec:solution}.

\subsection{Coherence}\label{sec:attr-coh}

The \emph{role} of a feature $u \in V'_{\text{mid}}$ is its pair of
weighted outgoing and incoming edge profiles, with each coordinate
scaled by the causal importance of its endpoint:
\begin{align}
(v_u^{\text{out}})_k &= W_{uk}\,\sqrt{\mathrm{Inf}_k} \nonumber \\
(v_u^{\text{in}})_k  &= W_{ku}\,\sqrt{\mathrm{Rel}_k} \nonumber \\
r(u) &= \bigl[\,v_u^{\text{out}}\,;\, v_u^{\text{in}}\,\bigr] \in \mathbb{R}^{2N}.
\label{eq:role}
\end{align}
Here $\mathrm{Inf}_k$ and $\mathrm{Rel}_k$ are the per-node influence
and relevance scores defined in Appendix~\ref{app:flow} via the Neumann
series on $G$. Outgoing coordinates are weighted by the
\emph{downstream} importance of their target ($\sqrt{\mathrm{Inf}_k}$),
incoming coordinates by the \emph{upstream} importance of their source
($\sqrt{\mathrm{Rel}_k}$), so the role vector concentrates on the edge
dimensions that actually carry causal mass. This is exactly the role
vector computed by every clusterer in our evaluation (the
\texttt{compute\_phi\_vectors} construction).

We commit to a single closed form for the coherence loss: the
within-supernode sum of squared Euclidean distances of each member's
role to the supernode centroid,
\begin{equation}
L_{\text{coh}}(\varphi) \;=\;
\sum_{S \in \pi(\varphi)}\; \sum_{u \in S}
\;\bigl\lVert r(u) - \bar r(S) \bigr\rVert_2^2,
\label{eq:Lcoh}
\end{equation}
where $\bar r(S) = \frac{1}{|S|} \sum_{u \in S} r(u)$ is the supernode
centroid in role space.
Eq.~\eqref{eq:Lcoh} is the $K$-means within-cluster sum of squares on
the Inf/Rel-weighted edge profiles. It is fixed once and reported
uniformly across all solvers: a spectral clusterer that internally
minimizes a normalized cut on a reweighted affinity, a modularity
clusterer that maximizes a degree-corrected score, and a direct
$K$-means clusterer all yield a partition $\varphi$, and each is
scored on the same $L_{\text{coh}}(\varphi)$ above. The internal
surrogate may differ; the reported number does not.

\paragraph{Why this choice.}
The Inf/Rel reweighting matches the role vector actually computed by
the methods we evaluate, and the squared-Euclidean form makes
$L_{\text{coh}}$ a well-defined function of the partition regardless of
which surrogate any solver used internally. The price is that
$\mathrm{Rel}$ inherits its dependence on the token-weight vector
$\tau$, currently set to SHAP values---the same fairness limitation we
flag for $L_{\text{cons}}$ (Section~\ref{sec:attr-cons}). A
steering-based replacement---defining $r(u)$ as the per-feature logit
perturbation $\Delta_u$ under intervention rather than the structural
edge profile, and reusing Eq.~\eqref{eq:Lcoh}---would directly measure
whether grouped features behave alike causally and is the natural
future iteration.

\subsection{Conservation}\label{sec:attr-cons}

Edges in an attribution graph carry signed linear contributions;
composing them along directed paths gives the causal mass that reaches
$V_{\text{logit}}$ from $V_{\text{emb}}$. Conservation is measured in
the same units. Working with magnitudes $|W|$, define the outgoing
transition matrix
\begin{equation}
P^{\text{out}}_{uv}
\;=\;
\frac{|W_{uv}|}{\sum_{w : (u, w) \in E'} |W_{uw}|}
\label{eq:Pout-reform}
\end{equation}
for each non-sink $u \in V'$, and the source vector $F^{\text{src}}$
equal to $\tau$ on $V_{\text{emb}}$ and zero elsewhere. The node-level
flow $\mathrm{Flow}(u) = [F^{\text{src}\top}(I - P^{\text{out}})^{-1}]_u$
is well-defined on a DAG and satisfies a per-node conservation law plus
the end-to-end identity $\sum_{\ell \in V_{\text{logit}}} \mathrm{Flow}(\ell)
= 1$ on the original (unpruned) graph. Aggregating to supernodes,
$\mathrm{Flow}^{SN}(S) = \sum_{u \in S} \mathrm{Flow}(u)$, similarly
balances at every interior supernode of $G_{SN}$. Full statements,
proofs, and the supernode lift are deferred to the flow note that
accompanies this paper.

A summary loses causal mass in two distinct ways.

\paragraph{Pruning loss.}
Edges and nodes removed in the cut to $(V', E')$ carry flow that no
longer reaches a logit. The fraction of source mass that survives
pruning is the \emph{flow completeness}
\begin{equation}
\mathcal{C}(V', E')
\;=\;
\sum_{\ell \in V_{\text{logit}}} \mathrm{Flow}_{V'}(\ell)
\;\in\; [0, 1],
\label{eq:Ccal}
\end{equation}
where $\mathrm{Flow}_{V'}$ denotes flow computed on the pruned graph
with $\tau$ restricted and renormalized. $\mathcal{C}(V', E')$ depends
only on $(V', E')$ and is partition-invariant.

\paragraph{Aggregation loss.}
Contracting features into supernodes destroys edge weight in two ways.
Edges between members of the same supernode become self-loops and are
dropped from $G_{SN}$. Edges between members of different supernodes
sum with sign in $W^{SN}_{ST} = \sum_{u \in S, v \in T} W_{uv}$, so
opposing-sign edges from $S$ to $T$ can partially or fully cancel.
Both forms of loss are visible directly on the edge magnitudes:
\begin{equation}
D_{\text{agg}}(\varphi)
\;=\;
1 \;-\;
\frac{\sum_{(S, T) \in E^{SN}} \lvert W^{SN}_{ST} \rvert}
     {\sum_{(u, v) \in E'} \lvert W_{uv} \rvert}.
\label{eq:Dagg}
\end{equation}
The numerator is the total magnitude of inter-supernode weight retained
in $G_{SN}$; the denominator is the total edge magnitude in the pruned
graph. $D_{\text{agg}}(\varphi) \in [0, 1]$, with $0$ iff every pruned
edge is preserved as a same-sign contribution to a distinct
inter-supernode block, and $1$ iff every pruned edge is absorbed (for
example $|\pi(\varphi)| = 1$, all mid-graph features in one supernode).

\paragraph{Conservation loss.}
Combining the two,
\begin{equation}
L_{\text{cons}}(V', E', \varphi)
\;=\;
\underbrace{\bigl(1 - \mathcal{C}(V', E')\bigr)}_{\text{prune loss}}
\;+\;
\underbrace{D_{\text{agg}}(\varphi)}_{\text{aggregation loss}}.
\label{eq:Lcons}
\end{equation}
The first term is the price of pruning and depends on $(V', E')$ and on
the source vector $\tau$; the second is the price of grouping and
depends only on $\varphi$ given $(V', E')$. No similarity matrix,
signature reweighting, or spectral relaxation appears in either
definition.

\paragraph{Fairness caveat on $\tau$.}
$D_{\text{agg}}$ is structural: it depends only on the signed edge
weights $W$ and the partition $\varphi$, with no role for $\tau$ or
$\ell$, so the aggregation half of $L_{\text{cons}}$ is fairness-clean. The
prune-loss half $1 - \mathcal{C}$ is not: it is precisely the
fraction of $\tau$-injected mass that fails to reach a logit, which
makes $\mathcal{C}(V', E')$ equal to the relevance-reachability of
$V_{\text{logit}}$. Our current pipeline sets $\tau$ to SHAP values,
so the prune loss inherits the bias of an external attribution
method. We treat $\tau$ as a free parameter of the problem
statement---any nonnegative distribution on $V_{\text{emb}}$ with
$\sum_e \tau_e = 1$ is admissible---and we flag a principled fairness
fix as open work: candidate alternatives include uniform $\tau =
1/|V_{\text{emb}}|$, a backward influence-side flow rooted at
$V_{\text{logit}}$ with logit weights $\ell$, and a two-sided
symmetrization combining forward and backward completeness.

\subsection{Complexity}\label{sec:attr-cplx}

A reader follows the summary by tracing causal paths from input tokens
to output logits. The relevant complexity is therefore the typical
distance from $V_{\text{emb}}$ to $V_{\text{logit}}$ in the supernode
graph, measured in supernode hops:
\begin{equation}
L_{\text{cplx}}(V', E', \varphi)
\;=\;
\frac{1}{|\mathcal{P}|}
\sum_{(e, \ell) \in \mathcal{P}}
d_{G_{SN}}(e, \ell),
\label{eq:Lcplx}
\end{equation}
where $\mathcal{P} = \{(e, \ell) \in V_{\text{emb}} \times V_{\text{logit}} : \ell \text{ reachable from } e \text{ in } G_{SN}\}$
and $d_{G_{SN}}(e, \ell)$ is the shortest-path length between $e$ and
$\ell$ in $G_{SN}$, counted in supernode edges. We set $L_{\text{cplx}}
= 0$ by convention when $\mathcal{P} = \emptyset$ (a degenerate
disconnected summary).

Eq.~\eqref{eq:Lcplx} is the plain average shortest-path length over
embedding-logit pairs; no flow weighting and no $\tau$ enters its
definition, so $L_{\text{cplx}}$ is fairness-clean in the sense of the
caveat in Section~\ref{sec:attr-cons}. Collapsing all mid-graph
features into a single supernode trivially hits the floor
$L_{\text{cplx}} = 2$ (one hop into the collapsed supernode, one hop
out), so unbounded grouping is checked by $L_{\text{cons}}$ through
$D_{\text{agg}}$ rather than rewarded here.

\section{Evaluation Methodology}\label{sec:eval}

A candidate solution to~\eqref{eq:Lsum} is a triple $(V', E', \varphi)$.
We score it along three axes: (i) the objective itself, decomposed by
term; (ii) the causal behaviour of the resulting summarization graph
under intervention; (iii) external utility for the human reader.
Decomposing in this way lets us compare solvers that minimize different
surrogates of the same problem on a common scale.

\subsection{Objective-aligned metrics}\label{sec:eval-obj}

Because $L$ is concrete (Section~\ref{sec:attr-instance}), every solver
is scored on exactly the same three numbers regardless of which
surrogate it minimized internally.
\begin{itemize}[leftmargin=*,topsep=2pt,itemsep=1pt]
\item \textbf{Coherence.} The realized $L_{\text{coh}}(\varphi)$ from
Eq.~\eqref{eq:Lcoh}: within-supernode sum of squared Euclidean
distances of the Inf/Rel-weighted role vectors to the supernode
centroid, identical to running \texttt{compute\_phi\_vectors} followed
by within-cluster inertia.
\item \textbf{Conservation.} $\mathcal{C}(V', E')$ and
$D_{\text{agg}}(\varphi)$ from Eqs.~\eqref{eq:Ccal}--\eqref{eq:Dagg},
reported separately. The first is partition-invariant given a fixed
pruner and acts as a sanity check on the pruning step; the second is
the per-partition quantity solvers compete on. $\mathcal{C}$ depends on
the choice of $\tau$, which defaults to SHAP values
(Section~\ref{sec:attr-cons}).
\item \textbf{Complexity.} $L_{\text{cplx}}(\varphi)$ from
Eq.~\eqref{eq:Lcplx}: average shortest-path length from
$V_{\text{emb}}$ to $V_{\text{logit}}$ in $G_{SN}$. We also report
$|\pi(\varphi)|$ as a descriptive statistic, though it is not part of
the objective.
\end{itemize}

\subsection{Causal-behaviour metrics}\label{sec:eval-causal}

$L_{\text{cons}}$ as defined in~\eqref{eq:Lcons} is intrinsic to the
graph. We additionally verify that the resulting $G_{SN}$ behaves
causally as the partition claims. The flow-based measures of the
companion flow note apply directly and are graph-intrinsic (no forward
passes through the model are required):
\begin{itemize}[leftmargin=*,topsep=2pt,itemsep=1pt]
\item \textbf{Flow faithfulness correlation}
$\mathrm{FaithCorr}_{\mathrm{Flow}}(\varphi)$: Spearman correlation
between supernode flow $\mathrm{Flow}^{SN}(S)$ and deletion-induced
flow drop $\Delta \mathcal{C}(S)$. A faithful partition ranks supernodes
by flow consistently with their causal contribution to logit-bound
flow.
\item \textbf{Flow-DAUC} and \textbf{Flow-IAUC}: areas under the
deletion and insertion curves over supernodes ranked by flow.
Lower DAUC and higher IAUC indicate that flow-ranked supernodes carry
the causal mass the rank suggests.
\item \textbf{Monotonicity} $\mathrm{Mono}(\varphi)$: fraction of
consecutive ranked supernodes whose deletion drops are non-increasing.
\item \textbf{Sufficiency} $\mathrm{Suff}_k^{\mathrm{Flow}}$ and
\textbf{comprehensiveness} $\mathrm{Comp}_k^{\mathrm{Flow}}$ at fixed
top-$k$: how much sink flow survives keeping or removing only the
top-$k$ supernodes.
\end{itemize}
Each of these scores a partition $\varphi$ using only $(V', E', \varphi)$
and the source vector $\tau$, with no reference to the solver that
produced $\varphi$.

\subsection{External utility metrics}\label{sec:eval-ext}

Two further metrics live outside the optimization problem but are the
practical reason summarization is useful.

\paragraph{Semantic coherence (LLM judge).}
For each supernode $S$, prompt a judge model with the auto-interpretation
of every member feature and ask for a coherence score in $[0, 1]$;
average over supernodes. This measures whether the partition aligns with
human-readable concepts, which $L_{\text{coh}}$ approximates structurally
but does not enforce.

\paragraph{Steering.}
For each supernode $S$ in a held-out evaluation set, scale the
activations of all member features by a common factor at inference time
and measure the resulting shift in the target logit. A partition
supports steering if intervening at the supernode level produces
predictable and concept-aligned logit shifts. This is the operational
test of whether $G_{SN}$ exposes causally controllable handles.

\subsection{Two-stage evaluation}\label{sec:eval-compare}

The pipeline factors into a pruning stage (producing $(V', E')$) and
an aggregation stage (producing $\varphi$), each with its own solver.
We evaluate the two stages independently. Pruning is scored through
the prune-loss component $1 - \mathcal{C}(V', E')$ alone, on a
candidate $(V', E')$ basis, and is reported separately. The protocol
below is the aggregation comparison, with the pruner held fixed.

\paragraph{Fixed pruner.}
A single pruning algorithm produces the $(V', E')$ that feeds every
aggregation solver in the comparison. $\mathcal{C}(V', E')$ is
therefore identical across rows, the prune-loss contribution to
$L_{\text{cons}}$ is the same constant in every row, and all
differences come from $D_{\text{agg}}$, $L_{\text{coh}}$, and
$L_{\text{cplx}}$.

\paragraph{Per-solver best-configuration reporting.}
Each aggregation solver has its own hyperparameter space; we sweep
exhaustively and report \emph{only the best configuration per solver}
in the comparison table. The hyperparameter spaces are:
\begin{itemize}[leftmargin=*,topsep=2pt,itemsep=1pt]
\item \textbf{ASG-spectral (ours)}: number of clusters $K$, mean
operator $F_{\text{mean}} \in \{\text{arith}, \text{harm}, \text{geo}\}$,
layer-decay rate $\lambda \in [0, 1]$.
\item \textbf{ASG-agglomerative (ours)}: same $K, F_{\text{mean}}, \lambda$,
with agglomerative linkage replacing the spectral relaxation.
\item \textbf{Greedy modularity}: target community count $K$ on the
symmetrized weighted adjacency.
\item \textbf{Spectral-RBF}: $K$, RBF kernel over the Inf/Rel-weighted
role vectors $r(u)$.
\item \textbf{$K$-means}: $K$ on the same role vectors.
\end{itemize}
``Best'' is decided by a single selection criterion fixed in advance
and shared across all solvers (the experiments section gives the exact
form). $K$ ranges over an eigengap-bounded candidate set common to
every solver. Different solvers may select different $K^\star$, which
is permitted: matching $|\pi(\varphi)|$ across solvers would force
sub-optimal configurations on methods whose natural granularity
differs, and the metrics already account for granularity via
$L_{\text{cplx}}$.

\paragraph{Reporting and dominance.}
For each solver's best configuration we report all metric families
from Sections~\ref{sec:eval-obj}--\ref{sec:eval-ext}. A solver
dominates if it is Pareto-better on the objective-aligned metrics
($L_{\text{coh}}, D_{\text{agg}}, L_{\text{cplx}}$) and at least
matches on causal-behaviour and external metrics. Because every
objective-aligned metric is intrinsic to the supernode graph---no
method-private similarity matrix or spectral surrogate enters its
definition---a solver that does well here is doing well at the stated
problem rather than at its own internal proxy.
.
% Inherits all macros from main.tex (\layer, \Ncut, \vol, \cut, \sign, \argmin,
% \argmax, definition / lemma / proposition theorem envs). Flow operators from
% measure.tex are written as \mathrm{...} to avoid coupling to whether the
% \DeclareMathOperator versions are present in the preamble.

\section{Problem Formulation}\label{sec:setup}

A summary of an attribution graph is determined by a pruned subgraph, a
partition of the surviving nodes into supernodes, and the supernode-level
adjacency that the pair induces. We state the problem in three parts.
Section~\ref{sec:input} fixes the input, Section~\ref{sec:output} fixes the
output and the induced summarization graph, and
Section~\ref{sec:objective} states the objective as the sum of three
losses---coherence, conservation, and complexity---together with
structural well-formedness constraints. Section~\ref{sec:attr-instance}
then commits to a concrete closed form for each loss term. We keep the
problem statement and its instantiation separable from any
\emph{solver}: $K$-means on stacked edge profiles, spectral clustering,
greedy modularity, and the block-coordinate strategy of
Section~\ref{sec:solution} may minimize different surrogates internally,
but every solver is scored on the same $L$ defined here.

\subsection{Input}\label{sec:input}

A directed acyclic attribution graph $G = (V, E, W)$ has vertex partition
$V = V_{\text{emb}} \cup V_{\text{mid}} \cup V_{\text{logit}}$ into
embedding, mid-graph, and logit nodes; edge set $E$; and signed weight
matrix $W \in \mathbb{R}^{N \times N}$ with $N = |V|$. We use the
convention that $W_{uv}$ is the weight on the directed edge $u \to v$,
so $W$ is sparse and strictly upper-triangular in any topological
ordering of $V$. Each node carries a layer index $\layer(\cdot)$;
embeddings sit at layer $0$ and logits at the final layer. Two side data
accompany $G$: a token-weight vector $\tau \in \mathbb{R}^{|V_{\text{emb}}|}$
on the source set (normalized SHAP values, $\sum_e \tau_e = 1$) and a
logit-weight vector $\ell \in \mathbb{R}^{|V_{\text{logit}}|}$ on the
sink set, identifying which inputs and outputs the summary should
explain.

\paragraph{Mid-graph nodes.}
In a cross-layer transcoder attribution graph, $V_{\text{mid}}$
originally contains both \emph{feature} nodes (interpretable
transcoder features) and \emph{reconstruction-error} nodes (the
residual that the transcoder fails to capture at each layer/position).
Error nodes are not interpretable as concept-bearing units, so we
restrict $V_{\text{mid}}$ to feature nodes throughout: all subsequent
references to ``members'' of a supernode, to ``role'' vectors, and to
the partition $\varphi$ apply only to feature nodes.

\subsection{Output}\label{sec:output}

A \emph{summary} is a triple $(V', E', \varphi)$:
\begin{itemize}[leftmargin=*,topsep=2pt,itemsep=1pt]
\item a pruned subgraph $V' \subseteq V$, $E' \subseteq E \cap (V' \times V')$,
with $V_{\text{emb}} \cup V_{\text{logit}} \subseteq V'$ forced;
\item a partition $\varphi : V' \to \pi(\varphi)$ assigning each surviving
node to exactly one \emph{supernode} $S \in \pi(\varphi)$, with every
$v \in V_{\text{emb}} \cup V_{\text{logit}}$ a singleton supernode.
\end{itemize}
The summary induces the \emph{summarization graph}
$G_{SN} = (\pi(\varphi),\, E^{SN},\, W^{SN})$ whose nodes are the
supernodes and whose adjacency aggregates the surviving weights:
\begin{equation}
W^{SN}_{ST} \;=\; \sum_{(u, v) \in E'} W_{uv}\,
\mathbf{1}[u \in S]\,\mathbf{1}[v \in T],
\label{eq:wsn}
\end{equation}
with $E^{SN} = \{(S, T) : W^{SN}_{ST} \ne 0,\; S \ne T\}$. The
summarization graph is the artifact the user reads; $(V', E', \varphi)$
are the decision variables that produce it.

\subsection{Objective}\label{sec:objective}

\paragraph{Three losses.}
A summary is good to the extent that (i) the members of each supernode
share a mechanistic role, (ii) $G_{SN}$ preserves the causal mass that
flows through $G$ from $V_{\text{emb}}$ to $V_{\text{logit}}$, and (iii)
the summary is small enough to read. We assemble these into a single
scalar objective:
\begin{equation}
\min_{(V', E', \varphi)}\;
L(V', E', \varphi)
\;=\;
L_{\text{coh}} \;+\; L_{\text{cons}} \;+\; L_{\text{cplx}}.
\label{eq:Lsum}
\end{equation}
Each term is named by what it measures, not by how it is measured:
\begin{description}[leftmargin=2.0em,style=nextline,topsep=2pt,itemsep=2pt]
\item[$L_{\text{coh}}(\varphi;\, V', E')$.] Within-supernode disagreement
on mechanistic role.
\item[$L_{\text{cons}}(V', E', \varphi;\, G)$.] Causal mass lost when
$G$ is replaced by $G_{SN}$.
\item[$L_{\text{cplx}}(V', E', \varphi)$.] Description length of the
summary.
\end{description}
The closed form of each term---and the multipliers that balance them---is
a property of a \emph{solver}, not of the problem.
Section~\ref{sec:attr-instance} fixes a concrete instantiation for
transcoder attribution graphs; Section~\ref{sec:eval} explains how
candidate solutions to~\eqref{eq:Lsum} are scored.

\paragraph{Structural constraints.}
Independent of the loss, a summary must be structurally well-formed:
\begin{description}[leftmargin=2.0em,style=nextline,topsep=2pt,itemsep=2pt]
\item[(F1) DAG integrity.] $G_{SN}$ is acyclic.
\item[(F2) Partition well-formedness.] $\varphi$ is a function (each
$v \in V'$ in exactly one supernode); embedding and logit nodes are
singletons.
\end{description}
(F1) and (F2) are properties any valid summary must satisfy, regardless of
which loss instantiation or solver is in play. Soft surrogates
(Lagrangian relaxations, penalty terms, dual variables) are introduced
only by specific solvers; the problem statement treats them as hard.

\section{Instantiation for the Transcoder Attribution Graph}\label{sec:attr-instance}

The losses of Section~\ref{sec:objective} are abstract. We now define each
one concretely for the transcoder attribution graph, where feature nodes
correspond to transcoder features and edge weights $W_{uv}$ encode the
linear causal contribution of $u$ to $v$ under the cross-layer
transcoder linearization. The point of this section is to pin down
\emph{what} is being measured; the survey of solvers and their surrogate
forms is the subject of Section~\ref{sec:solution}.

\subsection{Coherence}\label{sec:attr-coh}

The \emph{role} of a feature $u \in V'_{\text{mid}}$ is its pair of
weighted outgoing and incoming edge profiles, with each coordinate
scaled by the causal importance of its endpoint:
\begin{align}
(v_u^{\text{out}})_k &= W_{uk}\,\sqrt{\mathrm{Inf}_k} \nonumber \\
(v_u^{\text{in}})_k  &= W_{ku}\,\sqrt{\mathrm{Rel}_k} \nonumber \\
r(u) &= \bigl[\,v_u^{\text{out}}\,;\, v_u^{\text{in}}\,\bigr] \in \mathbb{R}^{2N}.
\label{eq:role}
\end{align}
Here $\mathrm{Inf}_k$ and $\mathrm{Rel}_k$ are the per-node influence
and relevance scores defined in Appendix~\ref{app:flow} via the Neumann
series on $G$. Outgoing coordinates are weighted by the
\emph{downstream} importance of their target ($\sqrt{\mathrm{Inf}_k}$),
incoming coordinates by the \emph{upstream} importance of their source
($\sqrt{\mathrm{Rel}_k}$), so the role vector concentrates on the edge
dimensions that actually carry causal mass. This is exactly the role
vector computed by every clusterer in our evaluation (the
\texttt{compute\_phi\_vectors} construction).

We commit to a single closed form for the coherence loss: the
within-supernode sum of squared Euclidean distances of each member's
role to the supernode centroid,
\begin{equation}
L_{\text{coh}}(\varphi) \;=\;
\sum_{S \in \pi(\varphi)}\; \sum_{u \in S}
\;\bigl\lVert r(u) - \bar r(S) \bigr\rVert_2^2,
\label{eq:Lcoh}
\end{equation}
where $\bar r(S) = \frac{1}{|S|} \sum_{u \in S} r(u)$ is the supernode
centroid in role space.
Eq.~\eqref{eq:Lcoh} is the $K$-means within-cluster sum of squares on
the Inf/Rel-weighted edge profiles. It is fixed once and reported
uniformly across all solvers: a spectral clusterer that internally
minimizes a normalized cut on a reweighted affinity, a modularity
clusterer that maximizes a degree-corrected score, and a direct
$K$-means clusterer all yield a partition $\varphi$, and each is
scored on the same $L_{\text{coh}}(\varphi)$ above. The internal
surrogate may differ; the reported number does not.

\paragraph{Why this choice.}
The Inf/Rel reweighting matches the role vector actually computed by
the methods we evaluate, and the squared-Euclidean form makes
$L_{\text{coh}}$ a well-defined function of the partition regardless of
which surrogate any solver used internally. The price is that
$\mathrm{Rel}$ inherits its dependence on the token-weight vector
$\tau$, currently set to SHAP values---the same fairness limitation we
flag for $L_{\text{cons}}$ (Section~\ref{sec:attr-cons}). A
steering-based replacement---defining $r(u)$ as the per-feature logit
perturbation $\Delta_u$ under intervention rather than the structural
edge profile, and reusing Eq.~\eqref{eq:Lcoh}---would directly measure
whether grouped features behave alike causally and is the natural
future iteration.

\subsection{Conservation}\label{sec:attr-cons}

Edges in an attribution graph carry signed linear contributions;
composing them along directed paths gives the causal mass that reaches
$V_{\text{logit}}$ from $V_{\text{emb}}$. Conservation is measured in
the same units. Working with magnitudes $|W|$, define the outgoing
transition matrix
\begin{equation}
P^{\text{out}}_{uv}
\;=\;
\frac{|W_{uv}|}{\sum_{w : (u, w) \in E'} |W_{uw}|}
\label{eq:Pout-reform}
\end{equation}
for each non-sink $u \in V'$, and the source vector $F^{\text{src}}$
equal to $\tau$ on $V_{\text{emb}}$ and zero elsewhere. The node-level
flow $\mathrm{Flow}(u) = [F^{\text{src}\top}(I - P^{\text{out}})^{-1}]_u$
is well-defined on a DAG and satisfies a per-node conservation law plus
the end-to-end identity $\sum_{\ell \in V_{\text{logit}}} \mathrm{Flow}(\ell)
= 1$ on the original (unpruned) graph. Aggregating to supernodes,
$\mathrm{Flow}^{SN}(S) = \sum_{u \in S} \mathrm{Flow}(u)$, similarly
balances at every interior supernode of $G_{SN}$. Full statements,
proofs, and the supernode lift are deferred to the flow note that
accompanies this paper.

A summary loses causal mass in two distinct ways.

\paragraph{Pruning loss.}
Edges and nodes removed in the cut to $(V', E')$ carry flow that no
longer reaches a logit. The fraction of source mass that survives
pruning is the \emph{flow completeness}
\begin{equation}
\mathcal{C}(V', E')
\;=\;
\sum_{\ell \in V_{\text{logit}}} \mathrm{Flow}_{V'}(\ell)
\;\in\; [0, 1],
\label{eq:Ccal}
\end{equation}
where $\mathrm{Flow}_{V'}$ denotes flow computed on the pruned graph
with $\tau$ restricted and renormalized. $\mathcal{C}(V', E')$ depends
only on $(V', E')$ and is partition-invariant.

\paragraph{Aggregation loss.}
Contracting features into supernodes destroys edge weight in two ways.
Edges between members of the same supernode become self-loops and are
dropped from $G_{SN}$. Edges between members of different supernodes
sum with sign in $W^{SN}_{ST} = \sum_{u \in S, v \in T} W_{uv}$, so
opposing-sign edges from $S$ to $T$ can partially or fully cancel.
Both forms of loss are visible directly on the edge magnitudes:
\begin{equation}
D_{\text{agg}}(\varphi)
\;=\;
1 \;-\;
\frac{\sum_{(S, T) \in E^{SN}} \lvert W^{SN}_{ST} \rvert}
     {\sum_{(u, v) \in E'} \lvert W_{uv} \rvert}.
\label{eq:Dagg}
\end{equation}
The numerator is the total magnitude of inter-supernode weight retained
in $G_{SN}$; the denominator is the total edge magnitude in the pruned
graph. $D_{\text{agg}}(\varphi) \in [0, 1]$, with $0$ iff every pruned
edge is preserved as a same-sign contribution to a distinct
inter-supernode block, and $1$ iff every pruned edge is absorbed (for
example $|\pi(\varphi)| = 1$, all mid-graph features in one supernode).

\paragraph{Conservation loss.}
Combining the two,
\begin{equation}
L_{\text{cons}}(V', E', \varphi)
\;=\;
\underbrace{\bigl(1 - \mathcal{C}(V', E')\bigr)}_{\text{prune loss}}
\;+\;
\underbrace{D_{\text{agg}}(\varphi)}_{\text{aggregation loss}}.
\label{eq:Lcons}
\end{equation}
The first term is the price of pruning and depends on $(V', E')$ and on
the source vector $\tau$; the second is the price of grouping and
depends only on $\varphi$ given $(V', E')$. No similarity matrix,
signature reweighting, or spectral relaxation appears in either
definition.

\paragraph{Fairness caveat on $\tau$.}
$D_{\text{agg}}$ is structural: it depends only on the signed edge
weights $W$ and the partition $\varphi$, with no role for $\tau$ or
$\ell$, so the aggregation half of $L_{\text{cons}}$ is fairness-clean. The
prune-loss half $1 - \mathcal{C}$ is not: it is precisely the
fraction of $\tau$-injected mass that fails to reach a logit, which
makes $\mathcal{C}(V', E')$ equal to the relevance-reachability of
$V_{\text{logit}}$. Our current pipeline sets $\tau$ to SHAP values,
so the prune loss inherits the bias of an external attribution
method. We treat $\tau$ as a free parameter of the problem
statement---any nonnegative distribution on $V_{\text{emb}}$ with
$\sum_e \tau_e = 1$ is admissible---and we flag a principled fairness
fix as open work: candidate alternatives include uniform $\tau =
1/|V_{\text{emb}}|$, a backward influence-side flow rooted at
$V_{\text{logit}}$ with logit weights $\ell$, and a two-sided
symmetrization combining forward and backward completeness.

\subsection{Complexity}\label{sec:attr-cplx}

A reader follows the summary by tracing causal paths from input tokens
to output logits. The relevant complexity is therefore the typical
distance from $V_{\text{emb}}$ to $V_{\text{logit}}$ in the supernode
graph, measured in supernode hops:
\begin{equation}
L_{\text{cplx}}(V', E', \varphi)
\;=\;
\frac{1}{|\mathcal{P}|}
\sum_{(e, \ell) \in \mathcal{P}}
d_{G_{SN}}(e, \ell),
\label{eq:Lcplx}
\end{equation}
where $\mathcal{P} = \{(e, \ell) \in V_{\text{emb}} \times V_{\text{logit}} : \ell \text{ reachable from } e \text{ in } G_{SN}\}$
and $d_{G_{SN}}(e, \ell)$ is the shortest-path length between $e$ and
$\ell$ in $G_{SN}$, counted in supernode edges. We set $L_{\text{cplx}}
= 0$ by convention when $\mathcal{P} = \emptyset$ (a degenerate
disconnected summary).

Eq.~\eqref{eq:Lcplx} is the plain average shortest-path length over
embedding-logit pairs; no flow weighting and no $\tau$ enters its
definition, so $L_{\text{cplx}}$ is fairness-clean in the sense of the
caveat in Section~\ref{sec:attr-cons}. Collapsing all mid-graph
features into a single supernode trivially hits the floor
$L_{\text{cplx}} = 2$ (one hop into the collapsed supernode, one hop
out), so unbounded grouping is checked by $L_{\text{cons}}$ through
$D_{\text{agg}}$ rather than rewarded here.

\section{Evaluation Methodology}\label{sec:eval}

A candidate solution to~\eqref{eq:Lsum} is a triple $(V', E', \varphi)$.
We score it along three axes: (i) the objective itself, decomposed by
term; (ii) the causal behaviour of the resulting summarization graph
under intervention; (iii) external utility for the human reader.
Decomposing in this way lets us compare solvers that minimize different
surrogates of the same problem on a common scale.

\subsection{Objective-aligned metrics}\label{sec:eval-obj}

Because $L$ is concrete (Section~\ref{sec:attr-instance}), every solver
is scored on exactly the same three numbers regardless of which
surrogate it minimized internally.
\begin{itemize}[leftmargin=*,topsep=2pt,itemsep=1pt]
\item \textbf{Coherence.} The realized $L_{\text{coh}}(\varphi)$ from
Eq.~\eqref{eq:Lcoh}: within-supernode sum of squared Euclidean
distances of the Inf/Rel-weighted role vectors to the supernode
centroid, identical to running \texttt{compute\_phi\_vectors} followed
by within-cluster inertia.
\item \textbf{Conservation.} $\mathcal{C}(V', E')$ and
$D_{\text{agg}}(\varphi)$ from Eqs.~\eqref{eq:Ccal}--\eqref{eq:Dagg},
reported separately. The first is partition-invariant given a fixed
pruner and acts as a sanity check on the pruning step; the second is
the per-partition quantity solvers compete on. $\mathcal{C}$ depends on
the choice of $\tau$, which defaults to SHAP values
(Section~\ref{sec:attr-cons}).
\item \textbf{Complexity.} $L_{\text{cplx}}(\varphi)$ from
Eq.~\eqref{eq:Lcplx}: average shortest-path length from
$V_{\text{emb}}$ to $V_{\text{logit}}$ in $G_{SN}$. We also report
$|\pi(\varphi)|$ as a descriptive statistic, though it is not part of
the objective.
\end{itemize}

\subsection{Causal-behaviour metrics}\label{sec:eval-causal}

$L_{\text{cons}}$ as defined in~\eqref{eq:Lcons} is intrinsic to the
graph. We additionally verify that the resulting $G_{SN}$ behaves
causally as the partition claims. The flow-based measures of the
companion flow note apply directly and are graph-intrinsic (no forward
passes through the model are required):
\begin{itemize}[leftmargin=*,topsep=2pt,itemsep=1pt]
\item \textbf{Flow faithfulness correlation}
$\mathrm{FaithCorr}_{\mathrm{Flow}}(\varphi)$: Spearman correlation
between supernode flow $\mathrm{Flow}^{SN}(S)$ and deletion-induced
flow drop $\Delta \mathcal{C}(S)$. A faithful partition ranks supernodes
by flow consistently with their causal contribution to logit-bound
flow.
\item \textbf{Flow-DAUC} and \textbf{Flow-IAUC}: areas under the
deletion and insertion curves over supernodes ranked by flow.
Lower DAUC and higher IAUC indicate that flow-ranked supernodes carry
the causal mass the rank suggests.
\item \textbf{Monotonicity} $\mathrm{Mono}(\varphi)$: fraction of
consecutive ranked supernodes whose deletion drops are non-increasing.
\item \textbf{Sufficiency} $\mathrm{Suff}_k^{\mathrm{Flow}}$ and
\textbf{comprehensiveness} $\mathrm{Comp}_k^{\mathrm{Flow}}$ at fixed
top-$k$: how much sink flow survives keeping or removing only the
top-$k$ supernodes.
\end{itemize}
Each of these scores a partition $\varphi$ using only $(V', E', \varphi)$
and the source vector $\tau$, with no reference to the solver that
produced $\varphi$.

\subsection{External utility metrics}\label{sec:eval-ext}

Two further metrics live outside the optimization problem but are the
practical reason summarization is useful.

\paragraph{Semantic coherence (LLM judge).}
For each supernode $S$, prompt a judge model with the auto-interpretation
of every member feature and ask for a coherence score in $[0, 1]$;
average over supernodes. This measures whether the partition aligns with
human-readable concepts, which $L_{\text{coh}}$ approximates structurally
but does not enforce.

\paragraph{Steering.}
For each supernode $S$ in a held-out evaluation set, scale the
activations of all member features by a common factor at inference time
and measure the resulting shift in the target logit. A partition
supports steering if intervening at the supernode level produces
predictable and concept-aligned logit shifts. This is the operational
test of whether $G_{SN}$ exposes causally controllable handles.

\subsection{Two-stage evaluation}\label{sec:eval-compare}

The pipeline factors into a pruning stage (producing $(V', E')$) and
an aggregation stage (producing $\varphi$), each with its own solver.
We evaluate the two stages independently. Pruning is scored through
the prune-loss component $1 - \mathcal{C}(V', E')$ alone, on a
candidate $(V', E')$ basis, and is reported separately. The protocol
below is the aggregation comparison, with the pruner held fixed.

\paragraph{Fixed pruner.}
A single pruning algorithm produces the $(V', E')$ that feeds every
aggregation solver in the comparison. $\mathcal{C}(V', E')$ is
therefore identical across rows, the prune-loss contribution to
$L_{\text{cons}}$ is the same constant in every row, and all
differences come from $D_{\text{agg}}$, $L_{\text{coh}}$, and
$L_{\text{cplx}}$.

\paragraph{Per-solver best-configuration reporting.}
Each aggregation solver has its own hyperparameter space; we sweep
exhaustively and report \emph{only the best configuration per solver}
in the comparison table. The hyperparameter spaces are:
\begin{itemize}[leftmargin=*,topsep=2pt,itemsep=1pt]
\item \textbf{ASG-spectral (ours)}: number of clusters $K$, mean
operator $F_{\text{mean}} \in \{\text{arith}, \text{harm}, \text{geo}\}$,
layer-decay rate $\lambda \in [0, 1]$.
\item \textbf{ASG-agglomerative (ours)}: same $K, F_{\text{mean}}, \lambda$,
with agglomerative linkage replacing the spectral relaxation.
\item \textbf{Greedy modularity}: target community count $K$ on the
symmetrized weighted adjacency.
\item \textbf{Spectral-RBF}: $K$, RBF kernel over the Inf/Rel-weighted
role vectors $r(u)$.
\item \textbf{$K$-means}: $K$ on the same role vectors.
\end{itemize}
``Best'' is decided by a single selection criterion fixed in advance
and shared across all solvers (the experiments section gives the exact
form). $K$ ranges over an eigengap-bounded candidate set common to
every solver. Different solvers may select different $K^\star$, which
is permitted: matching $|\pi(\varphi)|$ across solvers would force
sub-optimal configurations on methods whose natural granularity
differs, and the metrics already account for granularity via
$L_{\text{cplx}}$.

\paragraph{Reporting and dominance.}
For each solver's best configuration we report all metric families
from Sections~\ref{sec:eval-obj}--\ref{sec:eval-ext}. A solver
dominates if it is Pareto-better on the objective-aligned metrics
($L_{\text{coh}}, D_{\text{agg}}, L_{\text{cplx}}$) and at least
matches on causal-behaviour and external metrics. Because every
objective-aligned metric is intrinsic to the supernode graph---no
method-private similarity matrix or spectral surrogate enters its
definition---a solver that does well here is doing well at the stated
problem rather than at its own internal proxy.
